% Generated draft: requires mathematical and source review.
\documentclass[11pt]{article}
\usepackage[margin=25mm]{geometry}
\usepackage{fontspec}
\setmainfont{Latin Modern Roman}
\usepackage{amsmath,amssymb,mathtools,unicode-math}
\setmathfont{Latin Modern Math}
\usepackage{enumitem,needspace,xurl,hyperref}
\setlength{\parindent}{0pt}
\setlength{\parskip}{5pt}
\setlength{\emergencystretch}{3em}
\newcommand{\N}{\mathbb{N}}
\newcommand{\Z}{\mathbb{Z}}
\newcommand{\Q}{\mathbb{Q}}
\newcommand{\R}{\mathbb{R}}
\newcommand{\C}{\mathbb{C}}
\newcommand{\F}{\mathbb{F}}
\newcommand{\A}{\mathbb{A}}
\newcommand{\PP}{\mathbb{P}}
\newcommand{\E}{\mathbb{E}}
\newcommand{\eps}{\varepsilon}
\newcommand{\dd}{\,\mathrm{d}}
\begin{document}
\section*{UnsolvedMath ID: 800079}
\textbf{Source problem code:} catalogue entry
\subsection{Definitions and mathematical statement}
A group \(G\) is left-orderable if it admits a strict total order \(<\) such that \(g<h\) implies \(ag<ah\) for all \(a,g,h\in G\). A group is simple if it is nontrivial and its only normal subgroups are \(\{1\}\) and \(G\). It is finitely presented if it admits a presentation with finitely many generators and finitely many defining relations. It has type \(F_\infty\) if it admits a classifying CW complex \(K(G,1)\) whose \(n\)-skeleton is finite for every integer \(n\geq 0\).

For an abstract group \(G\), let \(H_b^n(G;\R)\) denote bounded group cohomology with trivial real coefficients: it is the cohomology of the subcomplex of bounded functions \(G^n\to\R\) in the usual inhomogeneous group cochain complex. With this coefficient convention, \(G\) is boundedly acyclic if
\[
H_b^n(G;\R)=0\qquad\text{for every integer }n\geq 1.
\]
The supplied question does not explicitly specify the coefficients in the term boundedly acyclic; the convention above follows its preceding discussion of real bounded cohomology.

\textbf{Existence questions.} Interpreting the source's word \emph{orderable} as \emph{left-orderable}, as suggested by its preceding discussion, ask:
\begin{enumerate}
\item Does there exist a finitely presented simple left-orderable group \(G\) such that \(H_b^n(G;\R)=0\) for every integer \(n\geq 1\)?
\item More strongly, does there exist a simple left-orderable group \(G\) of type \(F_\infty\) such that \(H_b^n(G;\R)=0\) for every integer \(n\geq 1\)?
\end{enumerate}
Type \(F_\infty\) implies finite presentability, so the second question strengthens the first. The source does not explicitly resolve whether \emph{orderable} means left-orderable or bi-orderable. If it means bi-orderable, the order must additionally satisfy \(g<h\Rightarrow ga<ha\) for all \(a,g,h\in G\), and that gives a different formulation. The left-orderable interpretation above requires confirmation from the source context.

The supplied background attributes to Navas [navas2018group, Question 8] the related question whether there exists a finitely generated left-orderable group \(G\) with no surjective homomorphism \(G\to\Z\) and with \(H_b^2(G;\R)=0\). It reports a positive answer by the authors named there as Fournier-Facioa and Lodha [fournier2023second]. It further reports that the simple orderable groups of Matte Bon and Triestino are boundedly acyclic [campagnolo2023algebraic, Corollary 1.16], whereas the only then-known finitely presented simple orderable groups are not boundedly acyclic [hyde2023finitely, Corollary 1.7]. It also cites two constructions obtained after dropping orderability [skipper2019simple; belk2022twisted], without providing their detailed hypotheses in this record. These are supplied contextual reports, not independently verified literature assessments. In particular, the absence of a surjection onto \(\Z\) and the degree-two vanishing condition belong to the related historical question; they do not replace the simplicity and all-positive-degree vanishing conditions of the present question.
\subsection{Short English statement}
The question asks whether a nontrivial simple group can simultaneously be finitely presented, admit a left-invariant total order, and have vanishing bounded cohomology with trivial real coefficients in every positive degree. It also asks for the stronger finiteness property \(F_\infty\). The meaning of \emph{orderable} and the bounded-cohomology coefficient convention require source confirmation. The registry labels the problem open; that label and the historical claims are supplied information, not independently verified conclusions.
\subsection{Sources}
\begin{itemize}
\item Local UnsolvedMath registry, permanent ID \(800079\), exact source problem code \texttt{catalogue entry}, title \emph{Foliations and Diffeomorphism Groups --- Question 20}. The supplied record attributes the problem to Nariman and Yazdi, proposed in 2025, and labels its status \emph{open}. Recorded upstream revision: \texttt{d371ee917ac9647fbb06cd6fc529deec6f3dda82}.
\item Nariman and Yazdi, \emph{Problem List on Foliations and Diffeomorphism Groups} (2025), Question 20. Supplied source URL: \url{https://arxiv.org/abs/2503.18274}. The record reports access on 2026-10-16 and extraction by \texttt{arxiv\_latex}. This access date and provenance are supplied metadata; no independent source visit or verification is asserted.
\item [skipper2019simple] and [belk2022twisted]: citation keys supplied for two constructions after dropping orderability. Full bibliographic entries and detailed construction statements are absent from the record.
\item Navas [navas2018group, Question 8]: supplied attribution of the related finitely generated left-orderable group question with no surjection onto \(\Z\) and vanishing \(H_b^2(G;\R)\). Full bibliographic details are absent.
\item [fournier2023second]: supplied attribution of a positive answer to that related question to authors written in the record as Fournier-Facioa and Lodha. Full bibliographic details are absent; the supplied spelling is retained.
\item [campagnolo2023algebraic, Corollary 1.16]: supplied reference for bounded acyclicity of the simple orderable groups of Matte Bon and Triestino. Full bibliographic details are absent.
\item [hyde2023finitely, Corollary 1.7]: supplied reference accompanying the claim that the only then-known finitely presented simple orderable groups are not boundedly acyclic. Full bibliographic details are absent. This is a source-specific historical qualification, not a new assessment of the literature.
\end{itemize}
\end{document}
